% TOP_PROBLEM: {"id": 53, "title": "Smale's 10th Problem: The Pugh Closing Lemma", "category_id": 6, "category": {"id": 6, "name": "geometry", "display_name": "Geometry", "description": "Euclidean and non-Euclidean geometry, geometric structures.", "slug": "geometry", "order_index": 6, "created_at": "2026-07-31T15:26:25.671Z"}, "status": "open"}
% FIRST_POSTED: 2026-09-27
% !TeX program = lualatex
% ATTEMPT_STATUS: unresolved
% Model: GPT-6 Astra Ultra; continuation dated 27 September 2026.
\documentclass[11pt]{article}
\usepackage[margin=25mm]{geometry}
\usepackage{fontspec}
\setmainfont{Latin Modern Roman}
\usepackage{amsmath,amssymb,mathtools}
\usepackage{unicode-math}
\setmathfont{Latin Modern Math}
\usepackage{xurl,hyperref}
\hypersetup{colorlinks=true,urlcolor=blue}
\setcounter{secnumdepth}{0}
\setlength{\parindent}{0pt}
\setlength{\parskip}{5pt}
\setlength{\emergencystretch}{3em}
\begin{document}
\section*{\texorpdfstring{Smale's 10th Problem: The Pugh Closing Lemma}{Smale's 10th Problem: The Pugh Closing Lemma}}\label{53}
\subsection{Definitions and mathematical statement}
Let $M$ be a compact smooth manifold and $\operatorname{Diff}^r(M)$ the space of $C^r$ diffeomorphisms with its $C^r$ topology. A point $p$ is nonwandering for $f$ if every neighborhood $V$ of $p$ satisfies $f^n(V)\cap V\ne\varnothing$ for some $n\ge1$. The closing assertion at differentiability order $r$ is
\[
 \forall f\ \forall p\in\Omega(f)\ \forall\mathcal U\ni f
 \quad\exists g\in\mathcal U\ \exists n\ge1:\ g^n(p)=p.
\]
Here $\mathcal U$ is a $C^r$ neighborhood and $\Omega(f)$ the nonwandering set. The question asks for the allowed differentiability orders of the source, in particular arbitrarily high finite orders; an infinite-order version uses the usual smooth topology. It concerns perturbations of diffeomorphisms without an added volume-preserving or symplectic restriction. Closing a point merely near $p$ must not replace periodicity of $p$ in this selected wording.
\par\smallskip\textit{Formulation and concept references:} \href{https://arxiv.org/html/2004.06855v2}{[S1]}.

\subsection{Short English statement}
Can an arbitrarily small perturbation in the required differentiability topology turn any nonwandering point of a diffeomorphism into a periodic point?
\subsection{Research attempt}
\paragraph{Exact point closing and source scope.}
The assertion for unrestricted higher-regularity diffeomorphisms
is \textbf{UNRESOLVED} here. The source [S1, Definition 1.1 and
Theorem A] explicitly makes the given point periodic and proves
that conclusion for partially hyperbolic maps with one-dimensional
center. Its introduction records the general $C^1$ theorem and
explains why some local perturbation techniques fail at higher
regularity. Neither a Hamiltonian-vector-field counterexample nor
a theorem restricted to a one-dimensional center settles the
unrestricted diffeomorphism question in the root statement.

We give a precise relocation lemma, explicit smooth subcases, and
a quantitative local closing construction. The latter makes the
derivative cost visible and identifies what recurrence alone fails
to estimate. Throughout, a return time $N$ need not be the least
period, since the target only asks $g^N(p)=p$.

\paragraph{Moving a nearby periodic point to the prescribed point.}
Fix a coordinate ball about $p$, a smaller concentric ball, and a
smooth bump function $\chi$ supported in the former and equal to
one on the latter. For $q$ close to $p$, the coordinate map
\[
 h_q(x)=x+\chi(x)(p-q)
 \tag{318.1}
\]
sends $q$ to $p$. Extend it by the identity outside the chart.
For $q$ sufficiently close, its derivative differs from the identity
by less than one in operator norm. The coordinate map is injective
because
$|h_q(x)-h_q(y)|\geq(1-\|Dh_q-I\|_\infty)|x-y|$.
It is a local diffeomorphism and proper, and agrees with the identity
outside a compact set, so is a global diffeomorphism. Equivalently,
one can use the associated proper degree-one local diffeomorphism.

All derivatives of $h_q-\mathrm{id}$ of each fixed order are
bounded by constants times $|p-q|$, since the bump support is fixed.
Thus $h_q\to\mathrm{id}$ in every finite $C^r$ topology and in
the smooth topology as $q\to p$. The inverses converge likewise
by the inverse derivative formulas.

If $q$ is periodic for a map $F$, put $G=h_qFh_q^{-1}$. Then
\[
 G^N(p)=h_qF^N(q)=h_q(q)=p.
 \tag{318.2}
\]
Composition and inversion are continuous in the stated
diffeomorphism topologies, so if $F\to f$ and $q\to p$ this
construction has $G\to f$. This proves an exact quantifier bridge:
arbitrarily close perturbations with periodic points arbitrarily
near $p$, with both approximations simultaneous, imply the selected
exact-point formulation. A theorem giving one periodic point
somewhere in $M$ does not satisfy that premise.

In particular, if $p\in\overline{\operatorname{Per}(f)}$, then
$p$ is $C^r$-closable without changing the conjugacy class of $f$.
The modification (318.1) is not claimed volume preserving; the
catalogue imposes no such restriction. Importing the same construction
unchanged into the conservative problem would require a different
relocation lemma.

\paragraph{Two complete classes at every differentiability order.}
For translations $R_\alpha(x)=x+\alpha$ of $\mathbb T^d$, choose
rational vectors $\alpha_j\to\alpha$. Each $R_{\alpha_j}$ has
a common finite period for all its points. The differences between
these maps have constant lifts, so their derivatives of every
positive order agree. Hence convergence is $C^\infty$, and every
point, not just a point near $p$, becomes periodic. If
$f=hR_\alpha h^{-1}$ for a fixed smooth diffeomorphism $h$, the
same argument with $hR_{\alpha_j}h^{-1}$ proves the conclusion.

For a hyperbolic toral automorphism induced by an integer matrix
$A\in GL(d,\mathbb Z)$, rational points are periodic. Indeed for
each denominator $q$, $A$ permutes the finite group
$(q^{-1}\mathbb Z/\mathbb Z)^d$, so every point in that group
has a finite orbit. Rational points are dense; (318.1)--(318.2)
therefore prove exact-point $C^\infty$ closing for every $p$.
Hyperbolicity is not even needed for this rational-grid observation,
although it gives another connection to the source's setting.

More generally, the standard Anosov closing/shadowing theorem
implies density of periodic points in the nonwandering set for an
Anosov diffeomorphism. Taking that theorem as an explicitly external
hyperbolic-dynamics input, the relocation lemma yields the selected
conclusion for such maps at their given regularity. This deduction
does not extend the shadowing theorem to arbitrary nonwandering
systems with a center direction.

\paragraph{A local construction with a fully quantified cost.}
Consider an orbit segment
\[
 x_0=p,\ x_i=f^i(p)\ (1\leq i\leq N),
 \qquad |x_N-p|=\delta
\]
in a coordinate chart near its two endpoints. Suppose a coordinate
ball $B(p,\rho)$ misses $x_1,\ldots,x_{N-1}$, and
$\delta<\rho/4$. Choose a fixed smooth bump $\chi_0$, equal to
one on the ball of radius $1/2$ and supported in the unit ball,
and define
\[
 h(x)=x+(p-x_N)\chi_0((x-p)/\rho).
 \tag{318.3}
\]
Then $h(x_N)=p$ and $h$ is the identity outside $B(p,\rho)$.
For a multi-index $\alpha$ with $|\alpha|=j$,
\[
 \|D^\alpha(h-\mathrm{id})\|_\infty
      \leq C_j\delta\rho^{-j}.
 \tag{318.4}
\]
If $C_1\delta/\rho<1$, the injectivity and properness argument
used above shows $h$ is a diffeomorphism. Set $g=h\circ f$.
Since $h$ fixes every intermediate image, induction gives
$g^i(p)=x_i$ for $0\leq i<N$ and $g^N(p)=p$.
For fixed $f$ and a fixed finite order $r$, the chain rule on a
compact manifold bounds
\[
 \|g-f\|_{C^r}\leq C(f,r)\max_{0\leq j\leq r}
                                      \delta\rho^{-j}.
 \tag{318.5}
\]
The coordinate norms are equivalent to any fixed finite atlas norms;
the constants incorporate that atlas and the chosen bump.

Consequently, any recurrent point admitting segments with
$\rho_j\to0$ and $\delta_j=o(\rho_j^r)$ under these separation
conditions is $C^r$-closable by this elementary method. If the ratios
vanish at every fixed derivative order along suitable segments, it
gives smooth closing as well. A basic smooth neighborhood constrains
only finitely many seminorms, so one chooses a segment adequate for
the largest required order; one need not control an infinite list
by one uniform bound.

\paragraph{Nonwandering is weaker than the recurrence used above.}
The definition of nonwandering gives points $q_j\to p$ and times
$N_j$ with $f^{N_j}(q_j)\to p$. It does not say that the orbit of
$p$ itself returns. If a construction closes these $q_j$ for maps
$F_j\to f$, then the relocation lemma closes $p$ exactly. This
shows how to use nonwandering returns without silently upgrading
them to recurrence of the specified point.

However, even for a recurrent point the endpoint distance $\delta$
alone says nothing about the radius of a ball avoiding all the
intermediate orbit points. Other visits can force $\rho$ to be
comparable to, or much smaller than, the endpoint error. The data
$\delta_j\to0$ and $\rho_j\to0$ do not imply
$\delta_j\rho_j^{-r}\to0$.

For a numerical scaling diagnostic, take $\rho_j=1/j$ and
$\delta_j=\rho_j^2$. The first derivative cost in (318.4) tends
to zero, while the second derivative cost stays of order one and
higher costs increase. These are possible relative scales, not a
constructed counterexample diffeomorphism. They refute the inference
from a small displacement to a small high-derivative perturbation.

\paragraph{Why smoothing a $C^1$ closing perturbation is insufficient.}
Suppose a $C^1$ theorem supplies a nearby map with a periodic point.
Approximating that map by smooth maps can preserve a hyperbolic
periodic orbit if hyperbolicity is available, but its distance from
the original $f$ in $C^r$ is still uncontrolled. Smoothness of an
approximant is not smallness of its high derivatives relative to $f$.
Formula (318.4) gives an explicit family of smooth modifications
exhibiting precisely this mismatch.

Nor can one perturb just the last return if the support intersects
earlier images: then $g^i(p)=f^i(p)$ may fail before time $N$,
invalidating the asserted closing identity. A valid multi-step
construction must track all altered intermediate points and divide
the correction so that every required derivative stays small. That
is substantive dynamical work, not a matter of choosing a smaller
bump after the fact.

\paragraph{Verification and outcome.}
The finite diagnostics enumerate rational torus orbits exactly on
finite grids, check the conjugacy identity for finite permutations,
and record exact rational costs $\delta\rho^{-r}$ in the displayed
scaling example. They do not simulate or claim shadowing for an
arbitrary diffeomorphism. The relocation and support arguments above
establish the actual exact-point claims.

We have proved the selected formulation for translations, their
smooth conjugates, toral automorphisms and, using the declared
shadowing input, Anosov maps. We also proved a quantitative sufficient
return condition. The missing general step is to exploit arbitrary
nonwandering returns while overcoming their support/derivative costs
in the prescribed $C^r$ topology. The unrestricted higher-order
closing assertion remains \textbf{UNRESOLVED} here.
\subsection{Sources}
\begin{itemize}
\item[S1]Shaobo Gan and Yi Shi, C\textasciicircum{}r-Closing lemma for partially hyperbolic diffeomorphisms with 1D-center bundle, Problem 1.
\url{https://arxiv.org/html/2004.06855v2}.
\textit{Formulation source.}
\end{itemize}
\end{document}
